% !TEX root = making_of.tex
\chapter{Primeros pasos}
\section{Sofware necesario}
En primer lugar, como se ha explicado ya en el guión de la práctica, vamos a emplear la placa STM32F411E, de la compañía ST, para controlar el motor. En la jerga de los programadores de microcontroladores, esta placa es conocida como la \emph{Black Pill}. Una de sus ventajas, es que tiene un excelente soporte de sofware y herramientas de desarrollo.

En esta sección vamos a describir como a asumir que estamos trabajando con un ordenador personal con sistema operativo Linux, en concreto vamos a emplear Ubuntu 24.04.4 LTS.
	\begin{verbatim}
	https://ubuntu.com/download/desktop
	\end{verbatim}


Algunas herramientas básicas necesarias son:
\begin{enumerate}
	\item Visual studio Code (VSC).	
Será el entorno de desarrollo que emplearemos para escribir, compilar y descargar el código que utilizará la placa ST. También servirá para desarrollar el código en Python que nos permitirá comunicarnos y controlar la placa desde el ordenador. Para instalarlo en Ubuntu, lo ideal es descargar directamente el paquete .dev, de la página web de VSC,
\begin{verbatim}
	https://code.visualstudio.com/thank-you?dv=linux64_deb
\end{verbatim}
para instalarlo ir a la línea de comandos:
\begin{verbatim}
sudo apt install ./Descargas/code_*.deb	
\end{verbatim}

Ver descripción detallada en el capítulo \ref{capvs}
	
	\item STM32CubeMX: Es la herramienta gráfica oficial de ST. Sirve para 	configurar los pines, el reloj (clock) y generar el código base en C.
	\begin{verbatim}
	https://www.st.com/en/development-tools/stm32cubemx.html
	\end{verbatim}

	\item STM32CubeCLT:
	Tiene todas las herramientas de ST, necesarias para compilar, enlazar etc.
	Para instalarlo
	Descargar de la página web de ST,
	\begin{verbatim}
	https://www.st.com/en/development-tools/stm32cubeclt.html
	\end{verbatim}
	el que pone Debian Linux installer.
	Descomprimir,
	\begin{verbatim}
		unzip st-stm32cubeclt_1.21.0_27995_20260219_1804_amd64.deb_bundle.sh.zip
	\end{verbatim}
	Dar Permisos
	\begin{verbatim}
		chmod +x st-stm32cubeclt_1.21.0_27995_20260219_1804_amd64.deb_bundle.sh
	\end{verbatim}
Lanzar el instalador
\begin{verbatim}
sudo ./st-stm32cubeclt_1.21.0_27995_20260219_1804_amd64.deb_bundle.sh	
\end{verbatim}
Presentará en el terminal la licencia de uso, ir hasta el final y aceptarla.
Debería instalarlo en:
\begin{verbatim}
	/opt/st/stm32cubeclt_1.21.0/
\end{verbatim}
Es toolset STM32CubeCLT incluye:

\begin{itemize}
	\item GNU C/C++ for Arm® toolchain tales como arm-none-abi-gcc(compilador), arm-none-abi-nm (symbol viewer) y otros.
	\item GDB debugger cliente y servidor (Depurador). Depurar el programa ejecutándolo en la propia placa
	\item STM32CubeProgrammer (STM32CubeProg) utility (Programador) Descargar el ejecutable desde el ordenador a la placa.
	\item System view descriptor files (.SVD) for the entire STM32 MCU portfolio.
    \item Map file associating STM32 MCUs and MCU development boards to the appropriate SVD 
\end{itemize}

	\item Cmake y Ninja: fundamentales para gestionar la construcción del proyecto: Librerías, dependencias, etc.
	
	\item ST-Link Drivers: permite la comunicación entre el ordenador y la placa de desarrollo a través de un puertos USB. Estas dos últimas herramientas se instalan directamente desde VSC. (Ver capitulo \ref{capvs})
	\item Python y los modulos de Python: serial, struct, csv, sys, time, collections, datetime, pyqtgraph.
	\item (Opcional) Si se va a hacer el análisis de los datos y los modelos del motor en python, entonces es comveniente instalar también los módulos: pandas, numpy, matplotlib.pyplot y control
	
Para instalar STM32CubeMx, lo ideal es obtenerlo de la página web de ST,

\noindent https://www.st.com/en/development-tools/stm32cubemx.html

Igualmente, para STM32CubeProgrammer,

\noindent
https://www.st.com/en/development-tools/stm32cubeprog.html

El resto de las herramientas, se pueden instalar directamente a través de Visualstudio Code, Mediante las extensiones para SMT32.

En el caso de Python, es bastante cómodo instalar directamente anaconda.

https://www.anaconda.com/downloadP

Varios de los módulos de Python necesarios vienen instalados por defecto con Anaconda y el resto son fácilmente instalables empleando \texttt{conda}.
 
\end{enumerate}
